\abs
\vspace*{-1.5cm}

\begin{resumen}
Resumen en espa\~nol. Escribir un buen resumen es una tarea dif\'{\i}cil que solo puede
realizarse tras haber finalizado la escritura del trabajo. Ha de resumir los resultados
de la investigaci\'on permitiendo al lector hacerse una idea del trabajo sin que necesite
leerlo por completo. Tras leer el resumen otros investigadores deben ser capaces de
decidir si vale la pena leer el trabajo. Debe proporcionar informaci\'on suficiente
sobre los objetivos y logros del trabajo, y tal vez sobre las t\'ecnicas y herramientas
utilizadas.

Evita en la medida de lo posible notaci\'on matem\'atica en el resumen. Si es
imprescindible que aparezca una referencia bibliogr\'afica ha de ponerse entera:
el resumen debe ser una pieza de texto independiente, pues puede aparecer publicada
en otros documentos, bases de datos, websites, etc. Ten en cuenta que el resumen es la
parte de tu trabajo que m\'as gente va a leer. Merece la pena dedicarle
la atenci\'on necesaria.
\palabrasclave{Palabra clave 1 -- Palabra clave 2 --Palabra clave 3 -- $\ldots$.}
\end{resumen}

\vfil

\begin{abstract}
Abstract in English.
A good abstract is
difficult to write and can only be completed after the full dissertation has been
written. It represents a brief summary of the results of the dissertation research. By
summarising the results of the research, it allows other people to get an idea of what
was accomplished without having to read through the whole dissertation. Other
scholars can read an abstract to decide if looking at the full work will be worthwhile.
The abstract should provide sufficient information about the results of the research
that reading the full dissertation is not necessary, although your markers will read the
full dissertation. (\dots)
\keywords{Keyword 1 -- Keyword 2 -- Keyword 3 -- $\ldots$.}

\end{abstract}
